% Part: first-order-logic
% Chapter: introduction
% Section: substitution

\documentclass[../../../include/open-logic-section]{subfiles}

\begin{document}

\olfileid{fol}{int}{sub}

\olsection{代入}

我们将通过一个例子来说明各个部分是如何联系在一起的，以及语法和语义的展开如何为后续更深入的研究奠定基础。我们的推导系统应该允许从 $\lforall[\Obj
v_0][\Atom{\Obj P}]{\Obj v_0}$ 推导出 $\Atom{\Obj P}{\Obj a}$。我们甚至可能想把它陈述为一条推理规则。然而，要做到这一点，我们必须能够以最一般的形式来陈述它：不仅针对 $\Obj P$、$\Obj a$ 和 $\Obj v_0$，而且针对任意公式~$!A$、项~$t$ 和变元~$x$。（回想一下，常项符号是项，但我们也会考虑由常项符号和函数符号构成的更复杂的项。）因此，我们希望能这样表述：“只要推导出了 $\lforall[x][!A(x)]$，就有根据推断出~$!A(t)$——即去掉 $\lforall[x]$ 并将~$x$ 替换为~$t$ 的结果。”但“将 $x$ 替换为~$t$”究竟意味着什么？$!A(x)$ 和~$!A(t)$ 之间有什么关系？这总是行得通的吗？

为了精确表达这一点，我们定义\emph{代入}运算。代入其实并不简单，因为我们不能只是把~$!A$ 中所有的~$x$ 统统替换为~$t$，而且也不是任何项~$t$ 都能代入任何变元~$x$。我们将再次使用归纳定义来处理这个问题。但一旦定义完毕，将推理规则规定为“从 $\lforall[x][!A(x)]$ 推断 $!A(t)$”就成了一种精确的定义。此外，我们将能够证明这是一条“好”的推理规则，其意义在于 $\lforall[x][!A(x)]$ 语义地后承~$!A(t)$。但要证明这一点，我们又必须去证某个乍看会让你想问“我们为什么要证这个？”的结论。$\lforall[x][!A(x)]$ 语义地后承~$!A(t)$ 依赖于这样一个事实：$\Sat{M}{!A(t)}$ 是否成立仅取决于项~$t$ 的值，也就是说，如果我们令 $m$ 为项~$t$ 在~$\Domain{M}$ 中所指称的元素，那么 $\Sat{M}{!A(t)}[s]$ 当且仅当 $\Sat{M}{!A(x)}[\Subst{s}{m}{x}]$。即使 $t$ 包含变元，这一点也成立，但在精确陈述该结论时我们必须谨慎。

\end{document}